% !TeX root = examen_bimestral01_filaB.tex
% !TeX program = pdflatex
% Fila B. Mismas preguntas y puntajes que la fila A.
% Orden de V/F respecto de la fila A: 2, 4, 5, 1, 3.
\documentclass[10pt,a5paper]{article}
\newif\ifmostrarSoluciones
\mostrarSolucionesfalse
\usepackage{estilo_apuntes}
\dosPorCarillatrue% Use \dosPorCarillatrue para dos copias por carilla.
\encabezadoprueba{Examen bimestral 1}{Fila B}

\begin{document}
\begin{center}
  {\color{Azul}\large\bfseries Examen del primer bimestre · Fila B}
\end{center}
\noindent\textbf{Nombre:} \hrulefill\quad\textbf{Paralelo:} \makebox[15mm]{\hrulefill}
\medskip

\noindent\textbf{Indicaciones}
\begin{itemize}
  \item Duración: 1 hora y 30 minutos.
  \item Justifique sus respuestas de forma breve y precisa. Indique dominios, restricciones e hipótesis de los teoremas utilizados. Se valora el razonamiento.
  \item No se permite el uso de dispositivos electrónicos. Cualquier intento de copia penalizará la evaluación con una calificación total de cero.
\end{itemize}

\noindent\textbf{Bloque I. Verdadero o falso justificado (5 puntos).}

\smallskip
\noindent Indique si cada una de las siguientes afirmaciones es \textbf{verdadera (V) o falsa (F)}. Si es verdadera, justifíquela mediante un argumento matemático; si es falsa, presente un contraejemplo o un cálculo que la refute. Cada afirmación vale 1 punto. Marcar únicamente V o F, sin justificación, no otorga puntaje. Las respuestas deben ser breves y precisas.
\begin{enumerate}
  \item Si una función tiene una asíntota horizontal cuando $x\to+\infty$, entonces está acotada en todo su dominio.
  \item Si $f$ es continua y estrictamente positiva en $[0,1]$, entonces existe un número $m>0$ tal que $f(x)\geq m$ para todo $x\in[0,1]$.
  \item En radianes, se cumple que $\displaystyle\lim_{x\to0}\frac{\sin(3x)}{\sin(2x)}=1$.
  \item Sean $f\colon A\to B$ y $g\colon B\to C$. Si $g\circ f$ es inyectiva, entonces $f$ es inyectiva.
  \item Si $f$ está definida cerca de $0$, salvo posiblemente en $0$, y satisface $|f(x)-2|\leq x^2$ allí, entonces $\lim_{x\to0}f(x)=2$.
\end{enumerate}

\medskip
\noindent\textbf{Bloque II. Desarrollo (5 puntos).}

\smallskip
\noindent Justifique los pasos esenciales; no es necesario extender el desarrollo más allá de lo solicitado.
\begin{enumerate}[resume]
  \item \textbf{(2,5 puntos)} Para $k\in\mathbb{R}$, se define en $(-\pi/2,\pi/2)$ la función
  \[
    G_k(x)=\begin{cases}
      \dfrac{\sqrt{1+4x^2}-1}{x\sin(2x)},&x\neq0,\\[2mm]
      k,&x=0.
    \end{cases}
  \]
  Determine el valor de $k$ para que $G_k$ sea continua en $0$, justificando el límite mediante herramientas algebraicas y trigonométricas.

  \item \textbf{(2,5 puntos)} Considere la función
  \[
    h(x)=x^2+\sin x,\qquad x\in\mathbb{R}.
  \]
  Se desea probar que $h$ alcanza un mínimo absoluto en $\mathbb{R}$, sin calcular su valor.
  
  Tenga presente que el teorema de máximo--mínimo (Weierstrass) no se aplica directamente a todo $\mathbb{R}$. Sin embargo, es posible demostrar la afirmación anterior siguiendo los pasos que se indican a continuación:
  \begin{enumerate}
    \item \textbf{(0,75 puntos)} Justifique que $h$ alcanza un mínimo $m$ en el intervalo cerrado $[-1,1]$. Usando el valor $h(0)$, explique por qué $m\leq0$.
    \item \textbf{(0,75 puntos)} Pruebe que, si $|x|>1$, entonces $h(x)>0$.
    \item \textbf{(1 punto)} Utilice los resultados de los pasos anteriores para demostrar que $h(x)\geq m$ para todo $x\in\mathbb{R}$ y que ese valor se alcanza. Concluya así la proposición inicial: $m$ es un mínimo absoluto de $h$ en todo $\mathbb{R}$.
  \end{enumerate}
\end{enumerate}

\ifmostrarSoluciones
\newpage
\begin{solucion}[del examen bimestral · fila B]
\begin{enumerate}
  \item \textbf{Falso.} La asíntota describe el comportamiento hacia $+\infty$, no en todo el dominio. Por ejemplo, $f(x)=1/x$ en $(0,\infty)$ tiene asíntota horizontal $y=0$ hacia $+\infty$, pero no está acotada superiormente porque $f(x)\to+\infty$ cuando $x\to0^+$.
  \item \textbf{Verdadero.} La continuidad en el intervalo cerrado y acotado $[0,1]$ garantiza que se alcanza un mínimo, digamos $m=f(x_m)$. Como todos los valores de $f$ son estrictamente positivos y $x_m$ pertenece al intervalo, $m>0$. Además, por ser el mínimo, $f(x)\geq m$ en todo $[0,1]$.
  \item \textbf{Falso.} Para $x\neq0$ suficientemente próximo a $0$,
  \[
    \frac{\sin(3x)}{\sin(2x)}
    =\frac32\frac{\sin(3x)}{3x}\frac{2x}{\sin(2x)}
    \longrightarrow\frac32.
  \]
  Que ambos argumentos tiendan a cero no elimina los factores de escala.

  \item \textbf{Verdadero.} Si $f(x_1)=f(x_2)$, al aplicar $g$ resulta $(g\circ f)(x_1)=(g\circ f)(x_2)$. Como la composición es inyectiva, se deduce $x_1=x_2$. Esta es la condición de inyectividad de $f$.
  \item \textbf{Verdadero.} La desigualdad equivale a
  \[
    2-x^2\leq f(x)\leq2+x^2.
  \]
  Ambas cotas tienden a $2$ cuando $x\to0$. El teorema del Sanduche garantiza el límite indicado. No hace falta que $f(0)$ esté definido.
  \item Para $x\neq0$ en el intervalo indicado, el denominador no se anula. Al racionalizar,
  \begin{align*}
    \frac{\sqrt{1+4x^2}-1}{x\sin(2x)}
    &=\frac{4x}{\sin(2x)(\sqrt{1+4x^2}+1)}\\
    &=\frac{2x}{\sin(2x)}\,
      \frac{2}{\sqrt{1+4x^2}+1}.
  \end{align*}
  El primer factor tiende a $1$ por el límite fundamental del seno, y el segundo tiende a $1$ por continuidad de la raíz y de las operaciones involucradas. El límite es $1$, así que la continuidad en $0$ exige y se obtiene con $k=1$.

  \item
  \begin{enumerate}
    \item La función es continua en $[-1,1]$, por ser suma de funciones continuas. En ese intervalo cerrado y acotado, máximo--mínimo garantiza un punto $c$ donde alcanza su mínimo $m=h(c)$. Como $0\in[-1,1]$, se cumple $m\leq h(0)=0$.
    \item Como $\sin x\geq-1$, se tiene $h(x)\geq x^2-1$. Si $|x|>1$, entonces $x^2>1$, de modo que
    \[
      h(x)\geq x^2-1>0.
    \]
    \item Dentro de $[-1,1]$, se cumple $h(x)\geq m$ por definición de mínimo. Fuera del intervalo, $h(x)>0\geq m$. Por tanto,
    \[
      h(x)\geq m=h(c)\qquad\text{para todo }x\in\mathbb{R}.
    \]
    Como $m$ se alcanza en $c\in[-1,1]$, es un mínimo absoluto en todo $\mathbb{R}$.
  \end{enumerate}
\end{enumerate}

\end{solucion}
\fi

\end{document}
